\chapter{Results and Discussion}
\label{chap:results}

\section{Where does the positive descriptive association come from?}

The exploratory comparison suggests that country-years governed by highly populist parties are more
corrupt. The model comparison shows why that pattern should not be interpreted as a within-country
effect. Pooled OLS gives a positive populism coefficient of 0.145 (95\% CI $[-0.024,0.315]$,
$p=0.093$). The between-country estimate is larger at 0.226, but also imprecise (95\% CI
$[-0.065,0.516]$, $p=0.127$). Once country effects are introduced, the sign flips: the country-FE
estimate is $-0.065$ and the two-way FE estimate is $-0.068$. Neither single-predictor within estimate
is significant at 5\%, but their signs are stable.

\begin{figure}[H]
  \centering
  \includegraphics[width=0.82\textwidth]{between_within.png}
  \caption{Between-versus-within comparison. Positive pooled and between-country point estimates
  become negative in country fixed-effects models; all four single-predictor estimates are imprecise
  after dependence-adjusted inference.}
  \label{fig:betweenwithin}
\end{figure}

The evidence therefore supports a careful formulation: countries with more populist governments tend
to have higher average corruption, but the estimate is uncertain and does not survive the shift to a
within-country comparison. Stable regional and institutional differences plausibly account for much
of the descriptive association.

\section{Composite corruption: controlled association}

\Cref{tab:twfe} presents the nested fixed-effects models. Adding log GDP per capita makes the negative
coefficient marginally more precise; adding electoral democracy leaves it close to $-0.064$. M2 and
M3 exclude zero at the 5\% level. The sign does not support the hypothesised positive same-year
relationship.

\begin{table}[H]
\centering
\caption{Two-way fixed-effects association models for composite corruption.}
\label{tab:twfe}
\begin{threeparttable}
\small
\begin{tabular}{llrrrr}
\toprule
Model & Added terms & Populism coefficient & 95\% CI & $p$ & Within-\Rtwo{} \\
\midrule
M1 & Country + year FE & -0.0677 & $[-0.139,\ 0.004]$ & 0.064 & 0.010 \\
M2 & + log GDP per capita & -0.0715 & $[-0.141,\ -0.002]$ & 0.043 & -0.059 \\
M3 & + electoral democracy & -0.0638 & $[-0.123,\ -0.004]$ & 0.036 & -0.196 \\
M4 & + lagged corruption & -0.0033 & $[-0.016,\ 0.010]$ & 0.615 & 0.885 \\
\bottomrule
\end{tabular}
\begin{tablenotes}[flushleft]\footnotesize
\item M1--M3 use the same 3,940 observations and 95 countries. M4 uses 3,845 observations and answers a
different short-run question after accounting for outcome persistence. Negative within-
\Rtwo{} can occur because the reported metric compares the transformed fit with a benchmark; it does
not reverse the coefficient interpretation.
\end{tablenotes}
\end{threeparttable}
\end{table}

\begin{figure}[H]
  \centering
  \includegraphics[width=0.75\textwidth]{association_stability.png}
  \caption{Stability of the populism coefficient across the static M1--M3 specifications. Country-
  clustered 95\% confidence intervals are shown.}
  \label{fig:stability}
\end{figure}

The coefficient is statistically detectable in the main controlled model but substantively small. A
one-within-country-standard-deviation increase in governing-party populism corresponds to an estimated
0.010 reduction in the 0--1 corruption index. This is approximately 3\% of the overall corruption
standard deviation, or 10\% of its within-country standard deviation. A typical election-to-election
change implies an even smaller movement.

Inference is stable under country-and-year clustering and Driscoll--Kraay standard errors. First
differences remain negative ($-0.024$) but are not distinguishable from zero ($p=0.195$). Alternative
government aggregations and originally proposed institutional controls also retain the negative sign
but reduce precision. The appropriate summary is therefore \emph{sign-robust but not
significance-robust}; it is not evidence that populism causes lower corruption.

\subsection{Heterogeneity by electoral democracy}

The populism-democracy interaction is statistically meaningful. The model-implied populism slope is
$-0.164$ at the 10th percentile of polyarchy and $-0.076$ at the median, with confidence intervals
excluding zero. At the 90th percentile it is $+0.027$ and imprecise. The negative same-year pattern is
therefore concentrated in less democratic contexts. This could represent genuine institutional
heterogeneity, but it could also reflect expert coding, regime trajectories, or conditioning on a
variable partly affected by populist government. We treat it as descriptive moderation rather than a
causal mechanism.

\subsection{Non-linear predictive benchmark}

The country-held-out XGBoost model obtains out-of-fold \Rtwo{} around 0.45, compared with 0.38 for a
linear model on the same folds. Nested tuning increases XGBoost performance to approximately 0.51 and
selects shallow, regularised trees. In both untuned and tuned fits, populism is the least important of
the three predictors; electoral democracy and GDP dominate SHAP importance.

\begin{figure}[H]
  \centering
  \includegraphics[width=0.72\textwidth]{shap_importance.png}
  \caption{Out-of-fold XGBoost feature importance. SHAP magnitude describes model use, not statistical
  significance or causality.}
  \label{fig:shap}
\end{figure}

This benchmark reinforces a limited conclusion: even a flexible model that is allowed to use
between-country information relies relatively little on populism to predict corruption.

\section{Temporal direction for composite corruption}

The simple two-lag populism-to-corruption model gives a negative lag-1 coefficient of $-0.048$
($p=0.031$) and a near-zero lag-2 coefficient. This is opposite to the hypothesised positive delayed
effect. In the reverse model, neither corruption lag is individually clear. Once the dependent
variable's own history is included, the primary Granger-style joint tests do not reject either null:

\begin{itemize}
  \item populism $\rightarrow$ corruption: joint $p=0.686$;
  \item corruption $\rightarrow$ populism: joint $p=0.241$.
\end{itemize}

\begin{figure}[H]
  \centering
  \includegraphics[width=0.76\textwidth]{causal_direction.png}
  \caption{Controlled lead-lag coefficients in both proposed directions. The simple lag-1
  populism-to-corruption estimate is negative, but the stricter joint Granger-style test is null.}
  \label{fig:direction}
\end{figure}

Placebo leads do not provide a strong warning: future corruption terms jointly have $p=0.841$ in the
reverse model, and future populism terms have $p=0.096$ in the forward model. First differences are
also null in both directions. AIC and BIC weakly favour one-year Granger windows, but longer-window
comparisons use progressively smaller samples and are not treated as model-selection evidence.

\subsection{Event-study evidence}

The populist-entry event study identifies 45 events. Pre-event coefficients are jointly null
($p=0.983$), which is consistent with no detectable pre-trend, but post-entry coefficients are also
jointly null ($p=0.228$). The estimates do not reveal a statistically detectable corruption increase
after entry.

\begin{figure}[H]
  \centering
  \includegraphics[width=0.78\textwidth]{populist_entry_event.png}
  \caption{Composite corruption around populist entry. Event time $-1$ is the reference. Neither the
  pre-trend nor post-entry joint test is significant.}
  \label{fig:event}
\end{figure}

The reverse corruption-shock event study finds a post-event joint signal, but governing populism is
already elevated before the shock and the pre-trend test is borderline. The sequence is therefore more
consistent with political episodes in which both variables cluster than with a clean corruption-first
mechanism.

Taken together, the directional evidence does not support a robust temporal effect for the composite
index. The small negative same-year association should not be converted into a claim about temporal
causation.

\section{Institutional decomposition}
\label{sec:decomposition}

The composite index may conceal relationships concentrated in one institution. \Cref{tab:decomp}
reports within-standardised M3 coefficients and primary Granger-style joint tests. Positive values are
oriented toward more corruption for every outcome, and FDR correction covers the four subtype
association tests and eight subtype-by-direction Granger tests.

\begin{table}[H]
\centering
\caption{Primary corruption-dimension results.}
\label{tab:decomp}
\begin{threeparttable}
\small
\begin{tabular}{lrrrrr}
\toprule
& \multicolumn{2}{c}{Contemporaneous M3} & \multicolumn{2}{c}{Granger joint $p$} \\
\cmidrule(lr){2-3}\cmidrule(lr){4-5}
Outcome & Std. $\beta$ & FDR $q$ & $P\rightarrow C$ & $C\rightarrow P$ & $N_{M3}$ \\
\midrule
Composite benchmark & -0.097 & -- & 0.686 & 0.241 & 3,940 \\
Executive & 0.021 & 0.685 & 0.496 & 0.870 & 3,940 \\
Public sector & -0.097 & 0.113 & 0.734 & 0.227 & 3,940 \\
Legislative & -0.076 & 0.213 & \textbf{0.003} & 0.064 & 3,745 \\
Judicial & \textbf{-0.173} & \textbf{0.009} & 0.316 & 0.251 & 3,940 \\
\bottomrule
\end{tabular}
\begin{tablenotes}[flushleft]\footnotesize
\item Granger cells show raw joint $p$ values for compactness. The legislative forward-direction test
has FDR $q=0.028$; all other subtype direction tests have $q>0.25$. The composite is a benchmark and is
excluded from subtype FDR families.
\end{tablenotes}
\end{threeparttable}
\end{table}

\subsection{Contemporaneous dimension results}

The negative composite association is not uniform. Executive corruption is near zero
($\widetilde{\beta}=0.021$, $q=0.685$). Public-sector and legislative corruption are negative but
imprecise after FDR correction. Judicial corruption has the clearest negative association
($\widetilde{\beta}=-0.173$, 95\% CI $[-0.284,-0.063]$, raw $p=0.002$, $q=0.009$).

\begin{figure}[H]
  \centering
  \includegraphics[width=0.82\textwidth]{decomposition_twfe.png}
  \caption{Controlled two-way fixed-effects association by corruption dimension. Effects are scaled by
  within-country standard deviations so that magnitudes are comparable.}
  \label{fig:decomposition}
\end{figure}

The common complete-case sample preserves the ranking. The judicial estimate remains negative
($\widetilde{\beta}=-0.144$, $q=0.042$), so the result is not explained by its larger available sample.
Its point estimate is unchanged under country, country-and-year, and Driscoll--Kraay covariance
estimators. A contemporaneous first-difference model retains the negative sign but does not survive
four-outcome correction (raw $p=0.033$, $q=0.134$). The coarse five-level judicial measure is also
negative ($\widetilde{\beta}=-0.137$, $p=0.026$). The sign is therefore robust to measurement and
inference choices, while confirmatory significance is weaker under differencing.

The judicial interaction mirrors the composite: the populism slope is negative at low and median
polyarchy and approximately zero at high polyarchy. This secondary result is consistent with
institutional heterogeneity, but it may also reflect low within-country resolution or expert coding in
less democratic regimes.

\subsection{Dimension-specific temporal ordering}

No subtype Granger test is significant after FDR except populism to legislative corruption (raw joint
$p=0.0035$, $q=0.0279$). The lag-1 coefficient is $-0.0328$ ($p=0.236$); the lag-2 coefficient is
$+0.0733$ (95\% CI $[0.0234,0.1231]$, $p=0.004$). In within-standardised terms, the two-year effect is
only $+0.0234$. The reverse legislative direction does not survive multiplicity correction
($p=0.064$, $q=0.257$).

\begin{figure}[H]
  \centering
  \includegraphics[width=0.74\textwidth]{decomposition_granger_fdr.png}
  \caption{FDR-adjusted joint Granger-style tests across corruption dimensions and directions. Only the
  populism-to-legislative-corruption family falls below 0.05.}
  \label{fig:decompgranger}
\end{figure}

The legislative signal passes the common-sample and placebo-lead checks: future populism does not
jointly predict current legislative corruption after FDR ($q=0.239$). However, it is not confirmed by
first differences (raw $p=0.374$), election-year observations ($p=0.701$), or government-change years
($p=0.486$). Simpler lead-lag models retain a signal in the full and forward-fill-restricted samples.

The evidence is consequently specific but fragile. A positive two-year coefficient is compatible with
slower legislative capture, but dependence on level-based, forward-filled information prevents a
causal interpretation. The decomposition refines the composite null result; it does not overturn it.


\section{Discussion}

Three findings matter. First, the eye-catching positive cross-country relationship is an imprecise
between-country pattern. It reverses sign in within-country models, indicating that stable differences
between countries are central. This is consistent with the very large regional and regime-type gaps
observed in the exploratory analysis.

Second, the composite panel does not reveal a robust temporal direction. Populism and corruption are
both persistent and may respond to the same political shocks. The event and placebo designs do not
support a general sequence in which populist entry is followed by rising corruption, nor one in which
corruption reliably brings populists to office. This absence of evidence describes what these annual
panel designs can detect; it is not proof that every causal effect is exactly zero.

Third, aggregation matters. Judicial corruption produces the strongest negative contemporaneous
association, while legislative corruption produces the only positive delayed signal. Neither pattern
justifies a simple causal story. The judicial result may capture coding and regime context, especially
because it is strongest at low democracy. The legislative result is theoretically compatible with
slower institutional capture \citep{chesterley2018,batory2016}, but its failure in first-difference and
election-only samples makes it provisional.

The results also caution against interpreting anti-corruption rhetoric as either proof of cleaner
government or proof of hidden corruption. Populism is a broad political style and ideology whose
institutional consequences vary with constraints, party organisation, and democratic context
\citep{guriev2022,huber2017}. Average panel effects can obscure that heterogeneity.

\section{Limitations}

\begin{enumerate}
  \item \textbf{Observational identification.} Fixed effects remove time-invariant confounders and
  year shocks, not unobserved country-specific changes. None of the models constitutes a natural
  experiment.
  \item \textbf{Measurement.} Both treatment and outcomes are expert-coded. Shared perceptions,
  retrospective coding, and reduced resolution in some countries can induce correlated error.
  \item \textbf{Forward-fill assumption.} Populism is held constant between elections. Annual sample
  size overstates the number of distinct political changes, and election-only estimates are less
  precise.
  \item \textbf{Persistence and dynamics.} Near-unit outcome persistence complicates level-based
  Granger tests. Dynamic fixed effects may be biased in finite panels, while first differences can be
  noisy.
  \item \textbf{Sample selection.} Strict coder filtering improves reliability but reduces coverage
  from 163 to 96 countries. Legislative missingness is linked to the existence of legislatures.
  \item \textbf{Temporal scope.} The analysis ends in 2019 and cannot evaluate the most recent wave of
  populist governments.
  \item \textbf{Multiplicity and dependence.} FDR correction limits false discoveries, but corruption
  dimensions are correlated and the legislative signal remains one result among several related
  specifications.
  \item \textbf{Average effects.} Country-level averages cannot identify heterogeneous effects by
  ideology, coalition structure, leader, or constitutional system without further design choices.
\end{enumerate}

\section{Conclusion}

This study asked whether governing-party populism is associated with political corruption, which
variable comes first, and which institution is most affected. A simple positive answer does not
survive the analytical sequence. The pooled and between-country coefficients are positive but
imprecise. The main within-country coefficient is small and negative, significant in controlled
static models but not under every robustness transformation. Composite lead-lag, Granger, event-study,
and first-difference evidence provides no robust temporal direction.

The institutional decomposition adds two qualifications. Judicial corruption shows the strongest
negative same-year association, particularly at lower democracy levels. Legislative corruption shows
a small positive two-year signal after the outcome's own history and multiplicity are considered, but
that result is not confirmed in first differences or election-only samples. The most defensible
conclusion is therefore a negative and nuanced one: the data do not support a broad claim that
populist government systematically raises political corruption, though a fragile delayed legislative
channel merits further research.

Future work should exploit quasi-experimental changes in governing coalitions, explicitly model
election timing and tenure, distinguish ideological families of populism, and combine expert indices
with administrative corruption outcomes. Those extensions would be better positioned to separate
institutional capture from perception, selection, and common political shocks.
